\documentclass[a4paper]{article}
\usepackage{amsfonts}
\usepackage{amsmath}
\usepackage{amssymb}
\usepackage{graphicx}     

\begin{document}
  
\noindent \large Auteur: Shahin Jamal \\
\noindent \large Studierichting: Informatica\\
\noindent \large Oefeningenreeks 2E nr. 9.6\\

\medskip

\normalsize

$4^{\tfrac{1}{x}} \cdot 16^{\tfrac{1}{x+2}} = 64^{\tfrac{1}{x+1}}$\\

$\Leftrightarrow 4^{\tfrac{1}{x}} \cdot 4^{\tfrac{2}{x+2}} = 4^{\tfrac{3}{x+1}}$\\

$\Leftrightarrow \dfrac{1}{x} + \dfrac{2}{x+2} - \dfrac{3}{x+1} = 0$\\

Zet alles op gelijke noemer.\\

$\Leftrightarrow \dfrac{(x+2)(x+1) + 2x(x+1) - 3x(x+2)}{x\cdot(x+2)(x+1)} = 0$\\

$\Leftrightarrow \dfrac{x^2+x+2x+2+2x^2+2x-3x^2-6x}{x\cdot(x+2)(x+1)} = 0$\\

$\Leftrightarrow \dfrac{-x+2}{x\cdot(x+2)(x+1)} = 0$\\

Enkel teller mag 0 zijn.\\

$\Rightarrow -x+2=0$\\

$x=2$\\

\begin{thebibliography}{999}
%\bibitem{a} Referentie
\end{thebibliography}
\end{document} 
